% Netzwerkanltung für die Studentenstadt Freimann
% Tex initially created by Maximilian Engelhardt <maximilian.engelhardt@stusta.mhn.de>

%\documentclass[a4paper,12pt,draft]{scrartcl}
\documentclass[a4paper,12pt]{scrartcl}

\usepackage[utf8]{inputenc}
%\usepackage{ngerman}
\usepackage{polski}
\usepackage{eurosym}
\usepackage{tabularx}
\usepackage[pdftex,final]{graphicx}
\usepackage{wrapfig}
\usepackage[top=1.5cm,bottom=2.5cm,left=1.5cm,right=1.5cm]{geometry}
%\usepackage[margin=2cm]{geometry}
\usepackage{hyperref}


\title{Miasteczko studenckie Studentenstadt Freimann:\\
       Instrukcja podłączenia do sieci komputerowej}
\date{\today}

\begin{document}

\maketitle

\begin{figure}[t!]
   \centering
   \vspace{-20pt}
   \includegraphics[width=0.8\textwidth,keepaspectratio]{Bilder/stunetp}
   \vspace{-20pt}
\end{figure}

\section*{Ogólne informacje nt. podłączania do sieci komputerowej}

Niniejszy dokument opisuje sposób podłączania komputera do sieci komputerowej w miasteczku studenckim Studentenstadt. Proszę wcześniej dokładnie przeczytać regulamin korzystania z sieci komputerowej (Benutzerordnung), który otrzymaliście Państwo przy podpisywaniu umowy wynajmu.

Do podłączenia się do sieci komputerowej potrzebny jest komputer wyposażony w gniazdo sieciowe typu Ethernet (obecne praktycznie we wszystkich obecnie produkowanych komputerach) oraz kabel sieciowy do podłączenia komputera do gniazdka sieciowego zamontowanego w pokoju. Kabel sieciowy można kupić w sklepach z elektroniką (np. Conrad), lub w czasie dyżuru StuStaNet e.V. (więcej na temat StuStaNet e.V. w dalszej części dokumentu).

Każdy, kto podłącza swój komputer do sieci komputerowej Studentenstadt jest odpowiedzialny za to, aby jego komputer nie stanowił zagrożenia dla pozostałych komputerów podłączonych do tejże sieci. W szczególności chodzi o zagrożenie wynikające z możliwości infekcji wirusem bądź innym szkodliwym programem. Dlatego w czasie korzystania z sieci komputerowej Studentenstadt należy pamiętać o:
\begin{itemize}
    \item regularnym instalowaniu wszystkich dostępnych aktualizacji systemu operacyjnego dot. bezpieczeństwa
    \item korzystaniu z regularnie aktualizowanego skanera antywirusowego
    \item korzystaniu z zapory sieciowej (Firewall)
\end{itemize}
W przypadku wykrycia przez oprogramowanie diagnostyczne sieci komputerowej Studentenstadt na podłączonym do niej komputerze obecności wirusa dostęp tego komputera do sieci zostanie natychmiast odcięty na czas usunięcia wirusa.

\begin{em}
Jeśli wyżej wymieniona sytuacja będzie się powtarzać dostęp do sieci komputerowej zostanie odcięty na stałe!!!
\end{em}

Tak drastyczne środki bezpieczeństwa są niestety konieczne, ponieważ wirusy wyrządzają ogromne szkody w sieci Studentenstadt. Leibniz-Rechenzentrum – operator łącza internetowego dla całego Studentenstadt w przypadku wykrycia w sieci wirusa blokuje dostęp do Internetu w całym miasteczku studenckim. 

\emph{Proszę mieć na uwadze:} Jeśli nie jesteście Państwo pewni, czy posiadacie wystarczającą wiedzę, aby móc zapewnić bezpieczeństwo na komputerze z zainstalowanym systemem MS Windows radzimy rozważyć skorzystanie z innego bezpiecznego systemu operacyjnego, np. Ubuntu Linux (dostępny bezpłatnie na stronie \url{http://www.ubuntu.com}). Systemy operacyjne z rodziny Linux oferują o wiele większe bezpieczeństwo, zbliżoną funkcjonalność oraz mogą być zainstalowane na danym komputerze równolegle z MS Windows bez konieczności rezygnacji z dotychczas używanego systemu.

\section*{Członkostwo w StuStaNet e.V.}

Członkostwo w StuStaNet e.V. Kosztuje jednorazowo \EUR{10} i przynosi wiele korzyści. Podstawową jest możliwość dostępu do Internetu bez konieczności korzystania z serwera proxy (czyli "normalnie").  Osoby nie będące członkami StuStaNet e.V., aby móc korzystać z Internetu muszą, w każdym programie korzystającym z Internetu (przeglądarka, Skype, itp.) odpowiednio skonfigurować serwer proxy (instrukcje jak to zrobić w przypadku popularnych przeglądarek znajdują się na końcu tego dokumentu), ponadto osoby takie będą miały zablokowany dostęp do niektórych usług internetowych (np. serwisy gier on-line jak kurnik.pl bądź streaming wideo). Ponadto członkowie StuStaNet e.V. mają do dyspozycji osobne konto poczty email, przestrzeń na serwerze Web z możliwością korzystania z PHP i bazy danych oraz wiele innych\footnote{\url{https://wiki.stusta.mhn.de/Liste\_der\_Vereinsdienste}}.

Wnioski o przyjęcie do StuStaNet e.V. można składać w czasie dyżurów StuStaNet e.V. w czasie semestru w poniedziałki i czwartki od 19.00 do 19.30 w Blaues Haus, na parterze w pokoju 0028.

Każdy kto jest zainteresowany funkcjonowaniem sieci komputerowej Studentenstadt i chciałby współuczestniczyć w jej rozwoju może dołączyć do zespołu adminów (admini są wybierani na początku każdego semestru osobno dla każdego budynku), lub przyjść na Adminrat, który ma miejsce w każdy pierwszy czwartek miesiąca o 20.00 w Teestube\footnote{Teestube znajduje się na 20 piętrze Hanns Seidel Haus (HSH), inaczej Grünes Haus, Christoph-Probst-Str. 16}.

\section*{Dalsze informacje}

Bardziej szczegółowe informacje dot. korzystania z sieci komputerowej oraz ogólnie życia w Studentenstadt można znaleźć na stronach Stusta-Wiki oraz na stronie info:

\begin{center}
  \begin{tabularx}{\linewidth}{|lX|}
    \hline
    \url{http://wiki.stusta.mhn.de/} & Wiki poświęcone życiu w Studentenstadt, można znaleźć informacje nie tylko nt. sieci komputerowej, ale też pobliskich sklepów, lekarzy, możliwości spędzania wolnego czasu i wiele innych.\\
    \hline
    \url{http://info.stusta.mhn.de/} & Ogłoszenia i ogólne informacje dot. życia w Studentenstadt\\
    \hline
  \end{tabularx}
\end{center}
W przypadku problemów z podłączeniem komputera do sieci można uzyskać pomoc u administratorów sieci Studentenstadt (Hausadministratoren). Lista administratorów z numerami pokojów i adresami email jest wywieszona w większości budynków na parterze. 

Administratorzy są odpowiedzialni wyłącznie za problemy związane (bezpośrednio, lub pośrednio) z dostępem do sieci.

\newpage
\section*{Wskazówki dot. zapewnienia bezpieczeństwa komputera.}

\subsection*{Programy antywirusowe}
\paragraph*{Windows}
\begin{center}
  \begin{tabularx}{\linewidth}{|p{.2\linewidth}XX|}
    \hline
    Nazwa & Strona internetowa & Uwagi\\
    \hline \hline
    Avira AntiVir & \url{http://www.free-av.com/} & Darmowy, na rynku od 1988 roku\\
    \hline
    Microsoft Security Essentials & \url{http://www.microsoft.com/security\_essentials/} & Darmowy program antywirusowy Microsoftu, wykrywa również oprogramowanie szpiegujące (spyware)\\
    \hline
    Sophos Antivirus & \url{http://sophos.lrz-muenchen.de/} & Komercyjne oprogramowanie antywirusowe, dostępne na stronie LRZ (link obok) za darmo dla studentów monachijskich uczelni\\
    \hline
    avast! & \url{http://www.avast.de/} & Darmowy dla użytkowników prywatnych\\
    \hline
  \end{tabularx}
\end{center}

Powyższa tabela zawiera jedynie wybrane programy antywirusowe, pełniejszą listę programów antywirusowych można znaleźć na stronach Wikipedii\footnote{\url{http://en.wikipedia.org/wiki/List\_of\_antivirus\_software}}, a ich recenzje na stronie AV-Comparatives.org\footnote{\url{http://www.av-comparatives.org/}}.

\paragraph*{Linux und MacOS X (Apple)}

Systemy operacyjne Linux oraz MacOS X nie wymagają używania programów antywirusowych.

\subsection*{Programy antyspyware}

Programy szpiegujące (spyware)  to programy, które bez wiedzy użytkownika gromadzą i wysyłają dane na jego temat (odwiedzane strony WWW, numery kart płatniczych, hasła)\footnote{\url{http://pl.wikipedia.org/wiki/Programy\_szpiegujące}}. Programy spyware nie stanowią bezpośrednio zagrożenia dla sieci, jednak zalecamy używanie oprogramowania zwalczającego również tego typu zagrożenia.

Przykładowe programy zwalczające spyware:

\begin{center}
  \begin{tabularx}{\linewidth}{|p{.18\linewidth}Xp{.3\linewidth}|}
    \hline
    Nazwa & Strona internetowa & Uwagi\\
    \hline \hline
    Windows Defender & \url{http://www.microsoft.com/germany/windows/products/winfamily/defender/default.mspx} & Wchodzi w skład Microsoft Security Essentials\\
    \hline
    Spybot – Search \& Destroy & \url{http://www.safer-networking.org/} & Darmowy dla użytkowników prywatnych\\
    \hline
    Ad-Aware & \url{http://www.lavasoft.com/products/select\_your\_product.php} & jw.\\
    \hline
  \end{tabularx}
\end{center}

\subsection*{Zapora sieciowa (Firewall)}

Zapora sieciowa kontroluje ruch sieciowy pomiędzy komputerem, a siecią, do której jest podłączony. Zadaniem zapory jest blokowanie i ostrzeganie użytkownika o ruchu sieciowym, który może być niebezpieczny (próby ataku, działalność wirusów i programów szpiegujących itp.). 

Każdy komputer podłączony do sieci Studentenstadt powinien mieć uruchomioną zaporę sieciową. Niemal każdy producent oprogramowania związanego z bezpieczeństwem oferuje również zapory sieciowe\footnote{\url{http://en.wikipedia.org/wiki/Comparison_of_firewalls}}, są w nie również standardowo wyposażone popularne systemy operacyjne (MS Windows, Linux). 


\newpage

\section*{Konfiguracja komputera}

Podłączenie komputera do sieci Studentenstadt wymaga wcześniejszego
odpowiedniego skonfigurowania go. Polega to głównie na wpisaniu w komputerze odpowiednich numerów IP. Każdemu gniazdku sieciowemu przypisane jest 8 kolejnych adresów IP. Zakres tych adresów znajduje się na etykiecie naklejonej na gniazdku sieciowym. W przypadku, gdy etykieta jest nieczytelna, lub zniszczona należy wypełnić dostępny w administracji formularz naprawy (Reparaturauftrag) i oddać go w tejże administracji (Hausverwaltung). Zniszczona etykieta zostanie wymieniona na nową.

Podłączenie komputera do sieci składa się z trzech kroków:

\begin{itemize}
    \item podłączenie komputera do gniazdka sieciowego
    \item konfiguracja ustawień sieciowych w systemie operacyjnym
    \item wprowadzenia ustawień serwera proxy (nie dotyczy członków StuStaNet e.V.)
\end{itemize}

\section*{Podłączenie do gniazdka sieciowego}

Komputer należy podłączyć do \emph{lewego} złącza.

UWAGA! \emph{W żadnym razie} nie należy do lewego złącza podłączać telefonów, modemów, urządzeń ISDN, lub tym podobnych. Może to doprowadzić do poważnego uszkodzenia sieci i urządzeń w niej pracujących.

\section*{Konfiguracja ustawień sieciowych:}

\begin{center}
  \begin{tabularx}{\linewidth}{|lXp{.2\linewidth}|}
    \hline
    Parametr & Wartość & Przykład \\
    \hline \hline
    Adres IP & \nolinkurl{10.150.xxx.yyy} -
    \nolinkurl{10.150.xxx.zzz}, \newline zakres 8 adresów IP podanych
    na gniazdku sieciowym, należy wybrać (jeden) dowolny z nich. & \nolinkurl{10.150.243.16} – \nolinkurl{10.150.243.23} \\
    \hline
    Maska podsieci & \nolinkurl{255.255.255.0} & \\
    \hline
    Brama sieciowa (ang. gateway) & \nolinkurl{10.150.xxx.254} \newline pierwsze trzy liczby jak w adresie IP, czwarta \nolinkurl{254} & \nolinkurl{10.150.243.254} \\
    \hline
    DNS-Serwer (Domain Name System) & \nolinkurl{10.150.127.2} \newline \nolinkurl{10.150.125.2} & \\
    \hline
    DNS-Suffix (nazwa domeny) & \nolinkurl{stusta.mhn.de} & \\
    \hline
    Skrypt proxy & \multicolumn{2}{l|}{\nolinkurl{http://wpad.stusta.mhn.de/proxy.pac}} \\
    \hline
    Serwer proxy & \multicolumn{2}{l|}{\nolinkurl{http://proxy.stusta.mhn.de:3128}} \\
    \hline
  \end{tabularx}
\end{center}

W przypadku ustawień proxy (ostatnie dwa wiersze w tabeli) należy podać albo skrypt, albo serwer. Zalecane jest podanie skryptu.

\subsection*{E-Mail, Skype, ICQ, gry online, ...}

Dostęp do sieci Internet na terenie Studentenstadt jest oferowany przez Studentenwerk München i ma służyć celom naukowym. Z tego powodu oraz ze względów bezpieczeństwa umożliwia on przeglądanie stron WWW za pośrednictwem serwera proxy, jednak dostęp do wielu innych usług jest zablokowany. Dotyczy to m.in. używania programów poczty elektronicznej, komunikatorów, gier online i in. Wymienione usługi dostępne są dla członków StuStaNet e.V.


\newpage
\enlargethispage{20pt}

\begin{figure}[t!]
    \raggedleft
    \vspace{-20pt}
    \includegraphics[height=1cm,keepaspectratio]{Bilder/Ubuntu_logo}
    \vspace{-20pt}
\end{figure}

\section*{Konfiguracja sieci krok po kroku}
\subsection*{Ubuntu Linux}
\begin{enumerate}
    \item Proszę otworzyć okno konfiguracji sieci System $\rightarrow$
      Einstellungen (Ustawienia)\footnote{Tłumaczenia niemieckich nazw
        menu mogą nie odpowiadać dokładnie nazwom użytym w polskiej
        wersji Ubuntu. W takim przypadku proszę o informację na listę
        admins@lists.stusta.mhn.de (po niemiecku).} $\rightarrow$ Netzwerkkonfiguration (Konfiguracja sieci).
    \item Proszę zaznaczyć w zakładce Kabelgebunden (połączenie za pomocą kabla) wpis odpowiadający Państwa karcie sieciowej (zazwyczaj "eth0") i kliknąć przycisk Bearbeiten (Edytuj).
    \item Proszę wybrać Ipv4-Einstellungen (Ustawienia IPv4) ustawić Methode (Metodę) na Manuell (Ręcznie).
    \item W Adressen (Adresy) proszę kliknąć na Hinzufügen (Dodaj).
      \begin{figure}[h!]
        \centering
        \begin{minipage}[c]{0.45\linewidth}
          \centering
          \includegraphics[width=\linewidth,keepaspectratio]{Bilder/IP_Gerneric}
        \end{minipage}
        \begin{minipage}[c]{0.5\linewidth}
          \centering
          \includegraphics[width=\linewidth,keepaspectratio]{Bilder/IP_Ubuntu}
          \caption{Przykładowa konfiguracja sieci w systemie Ubuntu}
          \vspace{-10pt}
        \end{minipage}
      \end{figure}
    \item Teraz należy podać adres IP, maskę podsieci, bramę sieciową, adres serwera DNS i domenę zgodnie z tabelą z sekcji Konfiguracja ustawień sieciowych i nacisnąć OK .
\end{enumerate}

\paragraph*{Ustawienie serwera proxy}
W Ubuntu istnieje możliwość globalnego ustawienia serwera proxy, tak aby nie musieć podawać jego danych w każdym programie osobno. W tym celu należy:
\begin{enumerate}
    \item Otworzyć Netzwerk-Proxy-Einstellungen(Ustawienia proxy) poprzez System $\rightarrow$ Einstellungen (Ustawienia) $\rightarrow$ Netzwerk-Proxy (Ustawienia proxy).
    \item Tutaj należy zaznaczać na samym dole opcję Automatische Proxy-Konfiguration(automatyczna konfiguracja proxy) w URL für Auto-Konfiguration(automatyczny skrypt konfiguracji): podać \url{http://wpad.stusta.mhn.de/proxy.pac}. Po czym zamknąć okno. 
\end{enumerate}
$\rightarrow$ Dostęp do Internetu jest w tej chwili prawidłowo skonfigurowany.



\newpage
\enlargethispage{20pt}

\begin{figure}[t!]
    \raggedleft
    \vspace{-20pt}
    \includegraphics[height=1cm,keepaspectratio]{Bilder/Windows_logo}
    \vspace{-20pt}
\end{figure}

\section*{Konfiguracja sieci krok po kroku}
\subsection*{MS Windows}
\begin{enumerate}
    \item Proszę otworzyć Panel Sterowania Start $\rightarrow$ Systemsteuerung(Panel Sterowania).
\end{enumerate}
\vspace{-15pt}
\paragraph*{Windows XP}
\begin{enumerate}
     \setcounter{enumi}{1}
     \item Wybrać Netzwerk- und Internetverbindungen(Połączenia sieciowe i internetowe).
     \item Następnie na samym dole okna wybrać Netzwerkverbindungen(Połączenia Sieciowe).
     \item Kliknąć prawym przyciskiem na LAN-Verbindung(Połączenie lokalne) i wybrać Einstellungen dieser Verbindung ändern(Właściwości).
\end{enumerate}
$\rightarrow$ Dalej jak w Windows Vista od pkt 5.
\vspace{-10pt}
\paragraph*{Windows Vista}
\begin{enumerate}
    \setcounter{enumi}{1}
    \item Proszę wybrać Netzwerk und Internet(Centrum sieci i odnajdywania) a następnie Netzwerkstatus und -aufgaben anzeigen(Zarządzaj połączeniami sieciowymi).
    \item Teraz należy wybrać Połączenie lokalne, a następnie Właściwości
    \setcounter{enumi}{4}
    \item W nowo otwartym oknie proszę zaznaczyć Internetprotokoll Version 4 (TCP/IPv4)(Protokół internetowy w wersji 4 (TCP/IPv4)) (Vista) ew. Internetprotokoll  (TCP/IP)(Protokół internetowy(TCP/IP)) (XP) i kliknąć Eigenschaften(Właściwości).
    \item Teraz należy podać adres IP, maskę podsieci, bramę sieciową i adres serwera DNS zgodnie z tabelą z sekcji Konfiguracja ustawień sieciowych.
      \begin{figure}[h!]
	\centering
        \vspace{-5pt}
        \begin{minipage}[c]{0.45\linewidth}
          \centering
          \includegraphics[width=\linewidth,keepaspectratio]{Bilder/IP_Gerneric}
        \end{minipage}
        \begin{minipage}[c]{0.48\linewidth}
          \centering
          \includegraphics[width=\linewidth,keepaspectratio]{Bilder/IP_Windows}
          \caption{Przykładowa konfiguracja sieci w systemie Windows (tutaj Vista, w XP identycznie)}
        \end{minipage}
      \vspace{-15pt}
      \end{figure}
    \item Następnie proszę kliknąć przycisk Erweitert(Zaawansowane) i w zakładce DNS na samej górze podajemy DNS-Suffix \nolinkurl{stusta.mhn.de}. Koniec. 
    \item Kliknąć OK.
\end{enumerate}




\newpage

\begin{figure}[t!]
    \raggedleft
    \vspace{-20pt}
    \includegraphics[height=1cm,keepaspectratio]{Bilder/OSXLeopard}
    \vspace{-20pt}
\end{figure}

\section*{Konfiguracja sieci krok po kroku}
\subsection*{Mac OS X}
\begin{enumerate}
    \item Proszę otworzyć okno konfiguracji sieci poprzez kliknięcie
      na Jabłko (lewy górny róg ekranu) i wybrać Systemeinstellungen(System preferences) $\rightarrow$ Netzwerk(Network).
    \item Proszę teraz wybrać urządzenie Ethernet.
    \item W zakładce TCP/IP ustawić opcję IPv4 Konfigurieren(Configure) na Manuell(Manually).
    \item Teraz należy podać adres IP, maskę podsieci, bramę sieciową i adres serwera DNS zgodnie z tabelą z sekcji Konfiguracja ustawień sieciowych.
      \begin{figure}[h!]
      \centering
%      \vspace{-5pt}
        \begin{minipage}[c]{0.38\linewidth}
          \centering
          \includegraphics[width=\linewidth,keepaspectratio]{Bilder/IP_Gerneric}
        \end{minipage}
        \begin{minipage}[c]{0.60\linewidth}
          \centering
          \includegraphics[width=\linewidth,keepaspectratio]{Bilder/IP_MAC}
          \caption{Przykładowe ustawienia sieci w Mac OS~X}
        \end{minipage}
%      \vspace{-15pt}
      \end{figure}
\end{enumerate}

\paragraph*{Ustawienie serwera proxy}
W Mac OS X istnieje możliwość globalnego ustawienia serwera proxy, tak aby nie musieć podawać jego danych w każdym programie osobno. W tym celu należy:
%\begin{wrapfigure}{r}{0.4\textwidth}
%  \vspace{-20pt}
%  \begin{center}
%    \includegraphics[width=0.48\textwidth,keepaspectratio]{Bilder/Proxy_MAC}
%  \end{center}
%  \vspace{-20pt}
%  \caption{Konfiguracja serwera proxy w Mac OS~X}
%  \vspace{-10pt}
%\end{wrapfigure}
\begin{enumerate}
    \item W poprzednim oknie kliknąć Weitere Optionen... (More Options) i wybrać zakładkę Proxies
    \item W panelu Zu konfigurierendes Protokoll (Protocols) wybrać Autom. Proxy-Konfiguration (Automatic proxy configuration) i jako URL podać \url{http://wpad.stusta.mhn.de/proxy.pac}. Następnie klikamy OK i w poprzednim oknie Anwenden(Apply). Można teraz zamknąć okno konfiguracji. Koniec.
\end{enumerate}
$\rightarrow$ Ustawienia sieciowe są w tej chwili skonfigurowane prawidłowo. 

\newpage

\section*{Ustawienia serwera proxy w przeglądarkach}

\begin{wrapfigure}{r}{0.5\textwidth}
%  \vspace{-20pt}
  \begin{center}
    \includegraphics[width=0.5\textwidth,keepaspectratio]{Bilder/Proxy_Firefox}
  \end{center}
%  \vspace{-20pt}
  \caption{Konfiguracja serwera proxy w przeglądarce Firefox}
%  \vspace{-10pt}
\end{wrapfigure}

\subsection*{\includegraphics[height=1.2cm,keepaspectratio]{Bilder/Firefox_35_logo} Mozilla Firefox}
\begin{enumerate}
    \item Proszę uruchomić Firefoxa
    \item W menu Extras (Narzędzia) wybrać Einstellungen (Opcje)
    \item Następnie wybrać kategorię Erweitert (Zaawansowane) i w niej zakładkę Netzwerk (Sieć).
    \item W sekcji Verbindung (Połączenie) kliknąć Einstellungen... (Ustawienia...).
    \item Pojawi się okno przedstawione po prawej stronie. Należy wybrać na samym dole Automatische Proxy-Konfigurations-URL (Adres URL automatycznej konfiguracji) i podać:\\ \url{http://wpad.stusta.mhn.de/proxy.pac}.
    \item Klikamy OK i zamykamy okno.
\end{enumerate}
$\rightarrow$ Przeglądarka Firefox jest teraz poprawie skonfigurowana.


\begin{wrapfigure}{r}{0.5\textwidth}
%  \vspace{-20pt}
  \begin{center}
    \includegraphics[width=0.5\textwidth,keepaspectratio]{Bilder/Proxy_IE}
  \end{center}
%  \vspace{-20pt}
  \caption{Konfiguracja serwera proxy w przeglądarce Internet Explorer}
%  \vspace{-10pt}
\end{wrapfigure}

\subsection*{\includegraphics[height=1.2cm,keepaspectratio]{Bilder/Internet_Explorer_7_Logo} Internet Explorer}
\begin{enumerate}
    \item Proszę uruchomić Internet Explorer.
    \item W menu Extras (Narzędzia) wybrać Internetoptionen (Opcje internetowe).
    \item W zakładce Verbindungen (Połączenia) kliknąć przycisk LAN-Einstellungen (Ustawienia sieci LAN).
    \item Zaznaczyć opcję  Automatisches Konfigurationsskript verwenden (Użyj skryptu automatycznej konfiguracji) i w polu Adresse (Adres) podać:\\ \url{http://wpad.stusta.mhn.de/proxy.pac}.
    \item Kliknąć OK i zamknąć pozostałe okna.
\end{enumerate}
$\rightarrow$ Przeglądarka Internet Explorer jest teraz poprawie skonfigurowana.

\newpage
\begin{wrapfigure}{r}{0.5\textwidth}
%  \vspace{-20pt}
  \begin{center}
    \includegraphics[width=0.5\textwidth,keepaspectratio]{Bilder/Proxy_Opera}
  \end{center}
%  \vspace{-20pt}
  \caption{Konfiguracja serwera proxy w przeglądarce Opera}
%  \vspace{-10pt}
\end{wrapfigure}

\subsection*{\includegraphics[height=1.2cm,keepaspectratio]{Bilder/Opera_O} Opera}
\begin{enumerate}
    \item Proszę uruchomić Operę.
    \item W menu Extras (Narzędzia) wybrać Einstellungen... (Preferencje...)
    \item Wybrać zakładkę Erweitert (Zaawansowane), w panelu po lewej stronie wybrać Netzwerk (Sieć) a następnie kliknąć przycisk Proxy (Serwery Proxy).
    \item W nowo otwartym oknie zaznaczyć ostatnią pozycję Automatische Proxykonfiguration verwenden (Skript) (Użyj automatycznej konfiguracji proxy) i w polu tekstowym wpisać:\\ \url{http://wpad.stusta.mhn.de/proxy.pac}.
    \item Zamknąć oba okna OK.
\end{enumerate}
$\rightarrow$ Przeglądarka Opera jest teraz poprawie skonfigurowana.


\end{document}

